% ============================================================================
% October 9, 2026 build (agent draft). Adapted from 04b_quantification.tex
% (Oct 6, base 15.07) to the 14.402 base point of the JMP deck of Oct 9.
% Fit, parameters and external inputs: latex/JMP_slides/prod_tables/
% (fit_full.tex, params.tex, external.tex; build_prod_tables.py, reading
% output/model/production/2007, promoted Oct 7: Mac round 3, chain 0, case
% 0010_nm, loss 14.402423836558775, selected_numerically_verified; receipt
% output/model/production/2007/PROMOTION_RECEIPT_20261007_mac_r3_chain0.md).
% Transition: H100 two-shock path, pair psi 0.1311 / 0.110
% (path_table_0.1311_0.110.md; figures
% output/model/jmp_slides_transition_figures_20261009/). 2023 table: deck frame
% "2023 Data and Model" (moments_2023_full_base14402.csv).
% ============================================================================
\providecommand{\mocktodo}[1]{\textcolor{red}{[TODO: #1]}}

\section{Quantification}
\label{sec:quant}

I treat the United States as an economy in transition to a smaller
population. The quantification has three steps. First, I calibrate a
stationary economy representing 2007, before the subsequent decline in
fertility, to match fertility, housing and wealth moments. Second, I hold
those parameters fixed and estimate a sequence of fertility shocks: two
unanticipated changes in the preference for children $\psi_t$ that match the
period fertility rate through 2023, carrying households and population
forward between the shocks. Third, I treat 2023 as a point on this
transition and project the economy forward, holding the preference for
children fixed after 2023. The policy experiments of
Section~\ref{sec:policy} all start from the same 2023 economy.

I organize the parameters into two groups. The first is set outside the
estimation, from external evidence or direct measurement. The second is
estimated by matching model moments to their empirical counterparts. I first
describe the data behind the targets and explain how annual data map into the
model's four-year period, because almost every parameter and moment passes
through that mapping.

\subsection{Data}
\label{sec:quant-data}

The targets combine five data sets, each chosen for the object it measures
best. Table~\ref{tab:quant-fit} lists every moment with its source. Fertility
moments come from the 2004 and 2006 CPS and the 2003--2006 natality records,
housing moments from the 2005--2006 ACS, the 2007 AHS and the PSID, and
wealth from the 2005 and 2007 PSID and the 2007 SCF.

\paragraph{Fertility.} The June Fertility Supplement of the CPS records the
number of children ever born to each woman. I pool the 2004 and 2006
supplements and measure childlessness and the share of mothers with exactly
one child among women aged 40 to 44, and the mean number of children ever born
among women aged 25. Children ever born are capped at three, the largest
number the model distinguishes. The timing of first births comes from the
NCHS natality files, which record every birth with the mother's age and the
birth order; I pool first births in 2003 to 2006.

\paragraph{Housing.} Mean rooms come from the 2007 AHS, for household heads
aged 18 to 85, without top-coding. Ownership among household heads aged 30 to
55 and the ownership gap between recent parents, whose oldest child is under
four, and heads without children come from the 2005--2006 ACS.
\mocktodo{state the dwelling-structure restriction used in both ACS ownership
moments (UNITSSTR codes 3--10 in the builder) and whether the model imposes it.}
The room response to the first birth is the PSID event-study estimate of
Section~\ref{sec:empirical-evidence} in the sample that fixes household
position before the birth: 1.465 rooms three to four years after the birth
(standard error 0.050). I use this sample rather than the household sample of
Figure~\ref{fig:emp-rooms-path}, which gives 1.03, because the model follows
the same households before and after the birth.

\paragraph{Wealth.} Aggregate net worth relative to aggregate annual gross
labor earnings comes from the 2005 and 2007 PSID: net worth of all household
heads aged 18 to 85, excluding business and farm equity and other real
estate, over gross labor earnings of heads aged 18 to 65. The bequest flow is
the annual flow of estates directed to children relative to aggregate net
worth in the 2007 Survey of Consumer Finances.

\subsection{From annual data to four-year periods}
\label{sec:quant-period}

A model period lasts four years. The unit of account is mean annual gross
labor earnings of working-age households, so a model value of one corresponds
to \$82,882 in the 2005 and 2007 PSID (2022 dollars). The mapping follows
three rules.

\paragraph{Stocks are measured in annual units; flows are four years long.}
Wealth, house prices and the bequest shifter $\theta_1$ are stocks and are not
rescaled. Income, consumption, rent and every other flow accrue over the whole
period, so a flow in the model is four times its annual counterpart. For
example, an entrant at age 18 with average efficiency earns
$4\times(1-0.080)\times0.721=2.65$ per period after the payroll tax, which is
0.72 times mean annual gross earnings. A moment that relates a stock to a flow
is computed against the annual flow, as in the data: the wealth target
divides net worth by annual gross earnings, and the bequest target divides the
estate flow over a period by four before dividing it by aggregate wealth.

\paragraph{Rates compound; hazards come from durations.}
Interest, discounting and depreciation compound over the four years. The gross
return is $R_b=1.02^4=1.082$, the period discount factor is
$\beta=\beta_{\text{ann}}^4$, and the period depreciation rate is
$1-(1-\delta_{\text{ann}})^4$. With $\beta_{\text{ann}}=0.940$, the period
discount factor is $0.781$. The property tax is the exception: it is a tax on
value levied each year, and I set the period rate to four times the annual
rate.%
\footnote{Compounding instead would give a period rate of 4.17 rather than
4.24 percent of value.}
Probabilities of events are period probabilities. Children leave home after
18 years on average, which I implement as a constant departure probability of
$\lambda=1/4.5=2/9$ per period for each child. Survival probabilities are
four-year probabilities between successive age cells. The probability $\pi_a$
that an attempt to have a child succeeds is the probability of a conception
within four years at that age.
\mocktodo{cite the source of the four-year conception probabilities and the
life table behind the survival probabilities.}
The earnings process is estimated directly at the four-year frequency rather
than converted from an annual process. I take non-overlapping four-year means
of household gross labor earnings in the annual PSID waves of 1984--1997,
remove age and year effects, and estimate $\rho_z=0.735$ from the first two
autocovariances of their logarithm. The implied annual persistence is
$0.735^{1/4}=0.926$.

\paragraph{Age-specific moments use the overlap of data ages with model cells.}
A model household of age $a$ occupies a four-year cell. A data statistic for
an age window weights each cell by the fraction of the cell that falls in the
window. Within a cell, a household's number of children is interpolated
linearly between its value at the start and at the end of the cell, which is
exact if births are spread uniformly over the four years. Women aged 40 to 44,
for example, overlap the cell $[38,42)$ for half its length and the cell
$[42,46)$ for three-quarters of its length, and children ever born at age 25
are read from the cell $[22,26)$, seven-eighths of the way from its start to
its end. For the mean age at first birth, I date a first birth in the model at
the midpoint of its cell and group the natality data into the same four-year
cells, so that both sides use the same representative ages.

The room response to the first birth is the one moment for which the period
length bounds how closely the model can mimic the data. In the model, I
compare households that have a first birth with otherwise identical
households that do not, one period, or four years, after the birth. The data
compare years three and four after the birth with years two and three before
it. The model's contrast is therefore a four-year response with no
pre-period, while the data's spans about six years.

\subsection{Parameters set outside the estimation}

Table~\ref{tab:quant-assigned} summarizes the parameters assigned outside the
estimation.

\paragraph{Demographics and preferences.}
Households enter at age 18, retire at 66, and live at most through age 85.
Survival is certain before retirement. Children remain at home for 18 years
on average. Births are possible through age 45, with the success probability
of an attempt falling from 0.98 at 18 to 0.50 at 42. I set risk aversion to
$\sigma=2$ and use the power equivalence scale of Scholz, Seshadri and Khitatrakun (2006) with
child weight $\omega=0.7$ and exponent $\nu=0.7$. The consumption weight
$\alpha_0=0.733$ is the expenditure share of nonhousing consumption among
childless renters in the 2019--2023 CEX. A birth into the state of three or
more children counts as 3.60 children when computing completed fertility: the
mean number of children, capped at five, among women aged 40 to 44 with three
or more in the 2024 CPS.

\paragraph{Housing finance and transaction costs.}
The annual real return on the bond is 2 percent, and selling a home costs 6
percent of its value (Greaney et al., 2025). Annual
depreciation is 1.4 percent (BEA and Federal Reserve, 2007) and the annual
property tax is 1.06 percent of value (ACS 2007--2011). The financed share is
$\phi=0.80$, so a buyer must hold at least 20 percent equity at the end of the
period in which it buys. There is no payment-to-income limit.

\paragraph{Income and pensions.}
The deterministic age profile of earnings is normalized to average one over
working ages. Efficiency follows the four-year process described above, with
$\rho_z=0.735$ and innovation standard deviation $\sigma_z=0.484$,
approximated with nine states. Entrants draw efficiency from its stationary
distribution. Every retiree receives the same pension. The payroll tax is 8.0
percent, and the pension is the one that balances the pension budget given
the age distribution of earnings: 22.9 percent of mean annual gross earnings,
which matches the ratio of pension income to gross earnings in the CPS ASEC.
% [CODE] e5f_stationary_paygo.py 47-62: the flat pension is solved from the
% payroll tax and the exogenous age-income mass (pension 0.9178 per period =
% 22.9% of mean annual earnings); e5f_social_security.py 50-96 certifies the
% balance.
\mocktodo{confirm the pension/earnings ratio's CPS ASEC year and definition.}

\paragraph{Housing products and supply.}
Owners choose among homes of $\{2,4,6,8,10\}$ rooms. Renters choose any size
up to $\bar h^R=6$ rooms, following Greaney et al.\ (2025).
The elasticity of housing supply with respect to the price is $\eta=0.63$,
from Baum-Snow and Han (2024).
\mocktodo{the deck cites Baum-Snow and Han (2024) for $\eta=0.63$; the
project notes record that no primary source for 0.63 has been confirmed
(supply\_elasticity\_history\_20261009.md). Check before submission; the
policy results of Section~\ref{sec:policy} depend more on the speed of
adjustment than on $\eta$.}

\paragraph{Entry wealth and bequests.}
Entrants' wealth follows the PSID rule of Section~\ref{sec:model}: net worth
of childless renters aged 18 to 24, censored at zero and assigned to
efficiency states by income rank. The bequest motive does not depend on the
number of children. I fix the bequest shifter at $\theta_1=0.008$, following
Greaney et al.\ (2025).

\begin{table}[t]
\centering
\caption{Parameters set outside the estimation}
\label{tab:quant-assigned}
\small
\begin{tabularx}{\textwidth}{@{}X l X@{}}
\toprule
Parameter & Value & Source or restriction \\
\midrule
Model period & 4 years & Model timing \\
Entry, retirement, last fertile ages & 18, 66, 45 & Life-cycle specification \\
Child departure per period $\lambda$ & $2/9$ & Mean stay of 18 years \\
Success of a birth attempt $\pi_a$ & 0.98 (18) to 0.50 (42) & Four-year conception probabilities \\
Risk aversion $\sigma$ & 2 & \\
Equivalence scale $(\omega,\nu)$ & $(0.7,\,0.7)$ & Scholz, Seshadri and Khitatrakun (2006) \\
Consumption share $\alpha_0$ & 0.733 & CEX 2019--2023, childless renters \\
Housing-supply elasticity $\eta$ & 0.63 & Baum-Snow and Han (2024) \\
Annual real interest rate & 2.0\% & Greaney et al.\ (2025) \\
Payroll tax rate $\tau^w$ & 8.0\% & CPS ASEC pension ratio \\
Financed share $\phi$ & 0.80 & 20\% down payment \\
Depreciation $\delta_H$; property tax $\tau^p$ & 1.4\%; 1.06\% per year & BEA and Fed 2007; ACS 2007--11 \\
Selling cost $\psi^s$ & 6\% & Greaney et al.\ (2025) \\
Owner house sizes $\mathcal H$ & 2, 4, 6, 8, 10 rooms & \\
Rental size cap $\bar h^R$ & 6 rooms & Greaney et al.\ (2025) \\
Bequest shifter $\theta_1$ & 0.008 & Greaney et al.\ (2025) \\
Earnings process $(\rho_z,\sigma_z)$ & 0.73, 0.48 per four years & PSID 1984--97, four-year earnings \\
Entry wealth & PSID levels & Childless renters 18--24, censored at zero \\
\bottomrule
\end{tabularx}
\end{table}

\subsection{Internally estimated parameters}

I choose eleven parameters jointly to match eleven moments, listed in
Table~\ref{tab:quant-fit}. Ten moments describe fertility, housing and wealth;
the eleventh is the normalization that completed fertility equal
replacement, 2.1, so that the 2007 economy is a stationary equilibrium with a
constant population.%
\footnote{In practice, for given values of the other ten parameters, the
house price solves the replacement condition and $H_0$ is the supply scale at
which that price clears the housing market with the population normalized to
one.}
Although all parameters are identified jointly, I organize the discussion by
economic block and describe which moments are most informative about each.

\paragraph{Fertility.}
Childlessness is most informative about the first-birth cost $\xi$. The mean
age at first birth and the number of children by age 25 discipline the timing
of first births, and with it the dispersion $\kappa_1$ of first-birth taste
shocks. The share of mothers with exactly one child is most informative about
the dispersion $\kappa_C$ for later births and the curvature $\gamma$ of the
child benefit. The level of the child benefit $\psi_0$ moves all fertility
moments together, and completed fertility pins it down given the price.

\paragraph{Housing and tenure.}
The room response around the first birth is the most informative moment for
the space requirement $h_P$: the larger the requirement, the more rooms a
household must add when its first child arrives. Mean rooms disciplines
housing demand more broadly, and with it the supply scale $H_0$. The
ownership rate and the ownership gap between recent parents and childless
households are most informative about the owner multiplier $\chi$ and the
dispersion $\kappa_T$ of tenure taste shocks.

\paragraph{Saving and bequests.}
The annual discount factor governs how much households carry across ages, and
aggregate net worth relative to aggregate annual gross earnings is the most
informative moment; the target is a ratio of totals, not an average of
household ratios. With $\theta_1$ fixed, the bequest flow relative to
aggregate wealth is most informative about $\theta_0$.

\begin{table}[t]
\centering
\caption{Internally estimated parameters}
\label{tab:quant-parameters}
\small
\begin{tabularx}{\textwidth}{@{}l X r X@{}}
\toprule
Parameter & Interpretation & Value & Most informative moments \\
\midrule
$\beta_{\text{ann}}$ & Annual discount factor & 0.940 & Wealth / earnings \\
$\xi$ & First-birth utility cost & 0.309 & Childlessness \\
$\kappa_1$ & First-birth taste scale & 0.119 & First-birth age; children by 25 \\
$\kappa_C$ & Later-birth taste scale & 0.358 & Exactly one child \\
$\psi_0$ & Initial preference for children & 0.156 & Completed fertility; all fertility moments \\
$\gamma$ & Curvature of the benefit of children & 0.098 & Exactly one child \\
$h_P$ & Parenthood space requirement (rooms) & 2.596 & First-birth room response \\
$\chi$ & Owner housing-service multiplier & 1.084 & Ownership; recent-parent gap \\
$\kappa_T$ & Tenure taste scale & 0.018 & Ownership; recent-parent gap \\
$H_0$ & Housing supply scale & 6.805 & Mean rooms; market clearing at replacement \\
$\theta_0$ & Bequest-motive scale & 0.228 & Bequests / wealth \\
\bottomrule
\end{tabularx}
\par\smallskip
\begin{minipage}{0.95\linewidth}
\footnotesize\textit{Notes:} Search bounds: $\beta_{\text{ann}}\in[0.93,0.99]$,
$h_P\in[0.1,2.6]$, $\kappa_1,\kappa_C\in[0.02,50]$. $H_0$ is reported on the
scale of \eqref{eq:m-supply}.
\mocktodo{$h_P=2.596$ is within 0.004 of its upper bound; decide how to report
this. The search stopped on disk space, so optimization convergence is not
certified (promotion receipt, Oct 7).}
\end{minipage}
\end{table}

\subsection{Target moments and fit}

Table~\ref{tab:quant-fit} reports the targets and the model's values. The
model matches completed fertility by construction, and matches childlessness,
the share with exactly one child, the mean age at first birth, ownership, the
recent-parent ownership gap and the wealth and bequest ratios closely. It
misses two moments, and the two misses account for most of the loss. Households
in the model have 0.55 children by age 25 against 0.81 in the data, so first
births start too late even though their mean age is matched. And the model's
family adds 1.26 rooms around the first birth against 1.47 in the data, with
the space requirement $h_P$ at its upper bound. Three moments that are not
targeted are close to the data: the share of first births at 30 or later, the
room gap between larger and smaller families, and the dispersion of wealth
among the old.

\begin{table}[t]
\centering
\caption{Targets and model moments, 2007 steady state}
\label{tab:quant-fit}
\small
\begin{tabularx}{\textwidth}{@{}X l r r@{}}
\toprule
Moment & Data & Target & Model \\
\midrule
\multicolumn{4}{@{}l}{\emph{Fertility}}\\
Completed fertility (normalization) & Replacement & 2.10 & 2.10 \\
Childless, women 40--44 (\%) & CPS 2004, 2006 & 19.8 & 20.2 \\
Exactly one child among mothers, 40--44 (\%) & CPS 2004, 2006 & 21.4 & 21.7 \\
Mean age at first birth (years) & NCHS 2003--2006 & 25.98 & 25.96 \\
Children ever born by age 25 & CPS 2004, 2006 & 0.810 & 0.553 \\
First births at age 30+ (\%)$^{\dagger}$ & NCHS 2003--2006 & 24.9 & 23.8 \\
\addlinespace
\multicolumn{4}{@{}l}{\emph{Housing}}\\
Mean occupied rooms & AHS 2007 & 5.73 & 5.82 \\
Ownership, heads 30--55 (\%) & ACS 2005--2006 & 67.6 & 67.0 \\
Recent-parent ownership gap & ACS 2005--2006 & 0.128 & 0.127 \\
First-birth room response & PSID event study & 1.465 & 1.264 \\
Rooms: 3+ versus 1--2 resident children$^{\dagger}$ & ACS 2005--2006 & 0.385 & 0.294 \\
\addlinespace
\multicolumn{4}{@{}l}{\emph{Wealth}}\\
Wealth / annual gross earnings & PSID 2005, 2007 & 4.46 & 4.58 \\
Annual bequests / wealth (\%) & SCF 2007 & 0.73 & 0.71 \\
Wealth/income p90 / median, ages 76--84$^{\dagger}$ & PSID & 3.52 & 3.31 \\
\bottomrule
\end{tabularx}
\par\smallskip
\begin{minipage}{0.95\linewidth}
\footnotesize\textit{Notes:} $^{\dagger}$ Not targeted. Samples and
definitions in Section~\ref{sec:quant-data}. Weighted loss 14.402, of which
the children ever born by 25 contribute 6.60 and the first-birth room
response 5.57.
\end{minipage}
\end{table}

\subsection{The fertility decline, 2007--2023}
\label{sec:quant-transition}

Figure~\ref{fig:quant-us-fertility} shows the motivation for the second step.
The period fertility rate of the United States peaked at about 2.1 in 2007
and fell to 1.62 by 2023. Completed fertility of women aged 40 to 44, which
reflects births over the previous two decades, was 1.92 in 2024. If period
fertility stays at its current level, completed fertility will converge to
1.62 within about two decades as the cohorts that had their children before
2007 age out.

\begin{figure}[t]
\centering
\includegraphics[width=\linewidth]{fertility_introduction_projection_1990}
\caption{Period and completed fertility in the United States}
\label{fig:quant-us-fertility}
\par\smallskip
\begin{minipage}{0.92\linewidth}
\footnotesize\textit{Notes:} Left: period total fertility rate, World
Development Indicators. Right: children ever born to women aged 40--44, CPS
Fertility Supplement. Dotted lines: a constant-rate scenario that holds period
fertility at its 2023 level; the completed-fertility continuation is rescaled
to end at that rate.
\end{minipage}
\end{figure}

Starting from the 2007 steady state, I hold all estimated parameters fixed
and let the preference for children change twice, each time unexpectedly and
believed to be permanent: from $\psi_0=0.156$ to $\psi_1=0.1311$ in 2007, and
to $\psi_2=0.110$ in 2015. The two values are chosen to match the period
fertility rate in the model's four-year periods 2011--2014 (1.861) and
2019--2022 (1.646). Between the shocks, households and the population are
carried forward, and prices, the pension and the rebate clear their markets
at every date. After 2023 the preference stays at $\psi_2$, households have
perfect foresight about prices, and the economy converges to a new stationary
equilibrium. Appendix~\ref{app:algorithms} describes the solution method.

Figure~\ref{fig:quant-transition-fit} compares the model's period fertility
with the data. The two shocks bring fertility to 1.86 in 2011--2014 and 1.65
in 2019--2022. After 2023, with the preference fixed, period fertility rises
slowly, to 1.79 by 2063. The recovery runs through the housing market: as the
smaller cohorts born after 2007 enter adulthood, the demand for rooms falls
and so does the price of space.

\begin{figure}[t]
\centering
\includegraphics[width=0.9\linewidth]{../../output/model/jmp_slides_transition_figures_20261009/fertility_fit_2007_2063.pdf}
\caption{Period fertility, model and data, 2007--2063}
\label{fig:quant-transition-fit}
\par\smallskip
\begin{minipage}{0.92\linewidth}
\footnotesize\textit{Notes:} Model: births per woman in each four-year period
of the two-shock transition from the 2007 steady state. Data: annual period
fertility and its four-year averages.
\end{minipage}
\end{figure}

Figure~\ref{fig:quant-full-transition} shows the whole transition, to 2403.
Period fertility returns to replacement after about a century and a half and
stays slightly above it for about a century before settling at replacement.
The number of households falls by about 40 percent,
to 0.594 of its 2007 level, total housing by about a quarter, and the house
price by 38 percent, to 0.485. Because the stock stays on the supply curve,
the fall in housing demand is split between a smaller stock and a lower price
in proportions set by $\eta$: with $\eta=0.63$, a stock 25 percent smaller
requires a price 38 percent lower. Along the way the price falls slowly. In 2023 it is 4.0 percent below
its 2007 steady-state level, and it is still falling in 2063.
% [MODEL] path_table_0.1311_0.110.md: benchmark price 0.7794; 2023 0.748;
% endpoint price 0.48467, population scale 0.59398, stock 4.312 vs 5.779.

\begin{figure}[t]
\centering
\includegraphics[width=\linewidth]{../../output/model/jmp_slides_transition_figures_20261009/full_transition_four_panel.pdf}
\caption{The full transition, 2007--2403}
\label{fig:quant-full-transition}
\par\smallskip
\begin{minipage}{0.92\linewidth}
\footnotesize\textit{Notes:} Period fertility, number of household heads,
total housing and the house price along the two-shock transition, with the
initial steady state normalized to 100. Diamonds mark the new stationary
equilibrium.
\end{minipage}
\end{figure}

The durability of the stock matters for this recovery.
Figure~\ref{fig:quant-full-transition-fixed} adds a path along which the
housing stock is frozen at its 2023 level. Because the rooms stay while the
population shrinks, the price falls faster and further before recovering
part of the fall, and period fertility rises above replacement from about
the 2070s, peaks near 2.2 around 2100, and settles back. The number of
households then stabilizes about 20 percent below its 2007 level rather than
40 percent. Space that cannot be withdrawn is passed on to smaller cohorts as
cheaper space, and cheaper space raises their fertility: completed fertility
of the cohort born in 2029 is 1.72 rather than 1.58, and of the cohort born
in 2057 1.93 rather than 1.72. Most of the gain is composition: fewer
households of these cohorts are still renting at ages 30 to 33, and renters
have the fewest children, although the renters who remain also gain.
% [MODEL] cohort_table.md (frozen stock vs stock that follows the price);
% "Born" in that table is the birth year of the cohort. Composition:
% notes/argument_for_corina_20261008.md step 17 (renters at 30-33 fall from
% a third to a fifth of a cohort; renters who remain 0.92 -> 1.01 for 2009).
\mocktodo{the cohort numbers come from the frozen-stock comparison of Oct 8
(cohort\_table.md); confirm they are on the same two-shock path as the
figure.}
Section~\ref{sec:policy-slow} introduces an intermediate case, a stock that
adjusts slowly toward the supply curve.

\begin{figure}[t]
\centering
\includegraphics[width=\linewidth]{../../output/model/jmp_slides_transition_figures_20261009/full_transition_four_panel_fixed.pdf}
\caption{The full transition with the housing stock fixed from 2023}
\label{fig:quant-full-transition-fixed}
\par\smallskip
\begin{minipage}{0.92\linewidth}
\footnotesize\textit{Notes:} As Figure~\ref{fig:quant-full-transition}, with
a second path (dashed) along which the housing stock is held at its 2023
level.
\end{minipage}
\end{figure}

\subsection{The 2023 economy}
\label{sec:quant-2023}

The 2023 economy is the starting point of every counterfactual, so I compare
it with the data on the same moments used in the calibration
(Table~\ref{tab:quant-2023}). None of these moments was targeted in the
transition, which matches only the period fertility rate. Housing is close to
the data: mean rooms are 5.54 against 5.58, and the room response at the first
birth and the room gap between larger and smaller families stay close to
their 2007 values in both the model and the data. Ownership among heads aged 30 to 55 is 69.6 percent against 58.7
percent: in the data ownership fell after 2007, while in the model, where
housing became cheaper, it rose. The wealth ratio is below the data, which
rose to 5.36 in 2017--2019.
\mocktodo{the 2023 wealth/earnings data value (5.36) is from the PSID 2017/19
on the narrow target definition, shown in red on the deck; confirm before
keeping it.}

\begin{table}[t]
\centering
\caption{The 2023 economy: data and model}
\label{tab:quant-2023}
\small
\begin{tabularx}{\textwidth}{@{}Xrr@{}}
\toprule
Moment & Data & Model \\
\midrule
Completed fertility, ages 40--44 & 1.92 & 1.71 \\
Childless, ages 40--44 (\%) & 18.8 & 27.4 \\
Exactly one child among mothers, 40--44 (\%) & 22.2 & 23.4 \\
Mean age at first birth (years) & 28.09 & 26.85 \\
First births at age 30+ (\%) & 37.7 & 28.2 \\
\addlinespace
Mean occupied rooms & 5.58 & 5.54 \\
Ownership, heads 30--55 (\%) & 58.7 & 69.6 \\
First-birth room response & 1.465 & 1.289 \\
Rooms: 3+ versus 1--2 resident children & 0.355 & 0.286 \\
\addlinespace
Wealth / annual gross earnings (PSID 2017/19) & 5.36 & 4.58 \\
Annual bequests / wealth (\%) & 0.88 & 0.80 \\
Wealth/income p90 / median, ages 76--84 & 3.45 & 3.31 \\
\bottomrule
\end{tabularx}
\par\smallskip
\begin{minipage}{0.95\linewidth}
\footnotesize\textit{Notes:} Model: the 2023 period of the two-shock
transition. Data: the sources of Table~\ref{tab:quant-fit} in their most
recent waves; wealth over earnings from the PSID 2017 and 2019.
\end{minipage}
\end{table}

The fertility rows are the model's largest miss, and they are informative
about what the preference shock can and cannot do. In the data, the decline
in period fertility since 2007 came with almost unchanged childlessness among
women aged 40 to 44 (19.8 percent in 2004--2006, 18.8 percent in 2024) and
with later first births: the mean age at first birth rose from 26.0 to 28.1
and the share of first births at 30 or later from 25 to 38 percent. In the
model, a lower $\psi_t$ works mainly on the extensive margin and on every
cohort alive. Women aged 40 to 44 in 2023 were 26 to 29 in 2007 and 34 to 37
in 2015, so both shocks lowered their first-birth hazard in their main
childbearing years, and 27.4 percent of them are childless. The data read the
decline as postponement among younger cohorts; the model reads it as a
permanent fall in the number of children of every cohort alive. A shock that
falls on young cohorts, or that raises the cost of early births, would bring
childlessness and the age at first birth closer to the data.
% [MODEL] childlessness_2023_miss_20261009.md.

The second miss concerns the house price. With only preference shocks,
falling fertility lowers the demand for family space, so the model has to
deliver a falling price: 4 percent below 2007 in 2023, and 10 percent below
with a slowly adjusting stock. In the data, the real house price rose by 15.6
percent between 2007 and 2023, and real rents by 15 to 19 percent. Part of the
decline in fertility since 2007 therefore came with a rising cost of space
that the estimation assigns to preferences. The transition says what housing
would have done in response to the fertility decline, not that housing caused
it.
% [DATA] Case-Shiller CSUSHPINSA / CPIAUCSL, 2023/2007 = 1.156
% (historical_supply_comparison.csv); CPI rent of primary residence, FRED.
\mocktodo{a refit with a housing-supply shock that targets the 2023 real price
as a third moment would turn this into a housing result
(housing\_tenure\_claims\_20261009.md, section 4).}

Figure~\ref{fig:quant-lifecycle-2023} compares life-cycle profiles in 2023.
The model reproduces the hump in housing size and the rise of ownership with
age, but its households buy too early, at the start of adult life, and keep
buying into old age: 40 percent of heads in their early twenties own in the
model against about 10 percent in the data, and almost all of the oldest heads
own against three quarters. Housing size peaks too early and falls too much
after 60, and children stay at home longer than in the data, because each
child leaves with a constant probability.
% [REVIEW Oct 10] lifecycle_2023.pdf and the deck do not state these numbers
% (40 percent, about 10 percent, three quarters); the figure file was not
% available to the review, so they are taken from the draft unverified.

\begin{figure}[t]
\centering
\includegraphics[width=\linewidth,trim=0 25 0 0,clip]{../../output/model/profile_2023_base14402/lifecycle_2023.pdf}
\caption{Life-cycle profiles in 2023, model and ACS}
\label{fig:quant-lifecycle-2023}
\par\smallskip
\begin{minipage}{0.92\linewidth}
\footnotesize\textit{Notes:} Homeownership, rooms and the share of households
with children at home by age of the household head. Model: the 2023 period of
the two-shock transition. Data: ACS 2023.
\end{minipage}
\end{figure}

Figure~\ref{fig:quant-allocation} returns to the allocation of large homes
documented in Section~\ref{sec:empirical-evidence}. The model reproduces the
shares held by young households without children and by middle-aged
households with children, but it gives young families with children twice
their share in the data (22.5 against 10.6 percent) and older households
without children half of theirs (21.6 against 40.0 percent). Part of the gap
is mechanical: in the model 9.3 percent of the large homes belong to old
households with a child at home, against 0.7 percent in the data, because a
child leaves with a constant probability. The model therefore understates the
intergenerational concentration of large homes. Any reallocation of space
from old to young that it predicts starts from a distribution in which the
young already hold more than they do in the data.
% [REVIEW Oct 10] 22.5 vs 10.6 and 21.6 vs 40.0 confirmed in
% notes/argument_for_corina_20261008.md (step 17 flag); 9.3 vs 0.7 not found
% in the notes and the figure file was not available to the review.

\begin{figure}[t]
\centering
\includegraphics[width=0.85\linewidth]{../../output/model/profile_2023_base14402/intergenerational_allocation_2023.pdf}
\caption{Who holds the large owner-occupied homes, model and data, 2023}
\label{fig:quant-allocation}
\par\smallskip
\begin{minipage}{0.92\linewidth}
\footnotesize\textit{Notes:} Share of owner-occupied homes with six or more
rooms held by each group, defined by the age of the householder and whether
a child lives in the home. Data: ACS 2023, as in
Figure~\ref{fig:large-homes}. Model: the 2023 period of the two-shock
transition.
\end{minipage}
\end{figure}
